\justifying
\NSFCBodyText

\subsubsubsection{研究内容}

本项目设置三项一一对应的研究内容，均在公开数据上完成，不依赖未授权采集。

\textbf{内容一：维吾尔语场景文字视觉编码。} 在 SUST 上训练或适配词级场景文字编码器，在 RUST 上测试。目标不是单独刷新识别排行，而是得到可供融合的文字区域特征 $s$。对无文字图像，令 $s=\mathbf{0}$ 并在融合中关闭该通道。中间指标为 RUST 识别准确率，用于确认编码器可用，不作为课题主验收。

\textbf{内容二：图像外观、场景文字与双语描述的异构融合。} 外观编码器 $f_v$、文字编码器 $f_s$ 与维/汉文本编码器分别提取 $v$、$s$、$h^{\mathrm{ug}}$、$h^{\mathrm{zh}}$。融合头 $g(v,s)$ 采用门控残差，避免文字噪声直接覆盖物体语义。对齐损失对外观—文本与文字—文本分开计算。比较简单拼接、单 InfoNCE 与门控双对齐，检验异构通道是否互相干扰。

\textbf{内容三：跨语言检索与小规模图像理解评价。} 在 Multi30k-Distant 测试集上报告图像到维/汉文本、维/汉文本到图像，以及汉查询—维描述、维查询—汉描述的 Recall@$K$ 与 MRR。在测试图像上拟建 200--400 题理解协议，区分“图中是什么”与“图中写了什么”。翻译流水线只作对照，不把 BLEU 当主指标。CUTE 与 MC$^2$ 仅用于文本编码器继续预训练。三项内容共享同一组基线与消融表，避免各写各的指标。经费与算力约束下，不把“再训练一个通用多语 VLM”列作研究内容。

\subsubsubsection{研究目标}

总体目标：给出可复现的维语场景文字增强多模态融合方法，并在公开基准上证明或证伪其对维—汉跨语言图像理解的贡献。

具体目标：（1）在 RUST 上得到可用的场景文字编码器，并使去掉该通道后的跨语言检索 Recall@5 下降可观测；（2）在 Multi30k-Distant 测试集上，完整模型的汉$\leftrightarrow$维互检索 Recall@5 高于无文字通道与 OCR—翻译流水线；（3）完成 200--400 题理解协议并报告准确率，不将其表述为已有大规模维语 VQA 库；（4）形成可支撑 1 篇中科院三区及以上 SCI 的方法与实验，署名遵守实验室第一单位要求。

\subsubsubsection{拟解决的关键科学问题}

\textbf{问题一：} 场景文字能否成为维语图像理解的稳定第三通道？若消融后检索与理解均不下降，则拒绝普遍有用的主张。

\textbf{问题二：} 外观、连写维文与维/汉描述如何融合而不互相干扰？若门控融合在物体检索和文字相关检索上同时差于单通道，则拒绝必须复杂融合的主张。

\textbf{问题三：} 共享图像作为枢纽，能否在缺少大规模维语 VQA 标注时改善汉$\leftrightarrow$维互检索？若翻译流水线全面不低于视觉枢纽，则把枢纽降为辅助模块。
